% B-linear-algebra
\chapter{Linear algebra of $2\times 2$ matrices}\label{app:linear-algebra}

This appendix proves the facts about real $2\times2$ matrices used in
\cref{ch:ode-toolbox,ch:turing}. Throughout, $B$ is a real $2\times2$ matrix, $I$
is the identity matrix, and $\abs{\cdot}$ is the Euclidean norm on $\R^2$ or
$\C^2$, $\abs{\mathbf z}^2=\abs{z_1}^2+\abs{z_2}^2$.

\section{Eigenvalues}\label{sec:linalg-eigenvalues}

\begin{proposition}\label{prop:2x2-eigenvalues}
Let $\Delta=\frac14(\tr B)^2-\det B$. The eigenvalues of $B$ are the roots of
$\mu^2-\tr B\,\mu+\det B=0$, namely
\begin{equation*}
  \mu_\pm=\tfrac12\tr B\pm\sqrt{\Delta},
  \qquad\text{where } \sqrt{\Delta}=i\sqrt{-\Delta} \text{ if } \Delta<0 .
\end{equation*}
If $\Delta>0$, they are real and distinct; if $\Delta=0$, there is one real
(double) eigenvalue $\frac12\tr B$; if $\Delta<0$, they are complex conjugate,
$\frac12\tr B\pm i\sqrt{-\Delta}$, with common real part $\frac12\tr B$. In all
cases $\mu_++\mu_-=\tr B$ and $\mu_+\mu_-=\det B$. A real eigenvalue has a real
eigenvector, and if $\bxi$ is an eigenvector for $\mu$, then $\bar\bxi$ is one for
$\bar\mu$.
\end{proposition}

\begin{proof}
Write $B=\left(\begin{smallmatrix}b_{11}&b_{12}\\b_{21}&b_{22}\end{smallmatrix}\right)$.
A number $\mu\in\C$ is an eigenvalue if and only if $B-\mu I$ is singular, i.e.\
$0=\det(B-\mu I)=(b_{11}-\mu)(b_{22}-\mu)-b_{12}b_{21}=\mu^2-\tr B\,\mu+\det B$.
The quadratic formula gives $\mu_\pm$ and the three cases, and comparing
coefficients in $\mu^2-\tr B\,\mu+\det B=(\mu-\mu_+)(\mu-\mu_-)$ gives the sum and
the product. If $\mu$ is real, $B-\mu I$ is a real matrix with determinant zero,
so its columns are linearly dependent over $\R$: $(B-\mu I)\bxi=0$ for some real
$\bxi\neq0$. Finally, conjugating $B\bxi=\mu\bxi$ gives $B\bar\bxi=\bar\mu\bar\bxi$,
since $B$ is real.
\end{proof}

\section{Norms}\label{sec:linalg-norms}

For a real or complex $2\times2$ matrix $M$ we use the \emph{operator norm}
$\norm{M}=\max\{\abs{M\mathbf z}:\mathbf z\in\C^2,\ \abs{\mathbf z}=1\}$, which
exists because $\mathbf z\mapsto\abs{M\mathbf z}$ is continuous on the compact unit
sphere. (In the main text, $\norm{\cdot}$ also denotes the $L^2$ norm of
functions; the argument always shows which is meant.) By definition
$\abs{M\mathbf z}\le\norm{M}\abs{\mathbf z}$, and
$\norm{\diag(c_1,c_2)}=\max\{\abs{c_1},\abs{c_2}\}$, because
$\abs{c_1z_1}^2+\abs{c_2z_2}^2\le\max_j\abs{c_j}^2\,\abs{\mathbf z}^2$ with equality
for a suitable unit vector. For real $\mathbf y\in\R^2$ and real $M$, the
Cauchy--Schwarz inequality gives
$\mathbf y\cdot M\mathbf y\le\abs{\mathbf y}\abs{M\mathbf y}\le\norm{M}\abs{\mathbf y}^2$,
and for $D=\diag(d_u,d_v)$ with $d_u,d_v>0$ we have
$\mathbf y\cdot D\mathbf y=d_uy_1^2+d_vy_2^2\ge\min\{d_u,d_v\}\abs{\mathbf y}^2$. Both
inequalities are used in the proof of \cref{thm:spectral-criterion}.

\section{The matrix exponential}\label{sec:linalg-exponential}

\begin{proposition}\label{prop:2x2-exponential-bound}
Let $\mu_\pm$ be the eigenvalues of $B$.
\begin{enumerate}[label=(\alph*)]
  \item Suppose $\re\mu_\pm<0$, and let $\gamma_0=-\max\{\re\mu_+,\re\mu_-\}>0$.
        Then for every $\gamma<\gamma_0$ there is $C\ge1$ such that
        \begin{equation*}
          \abs{e^{tB}\mathbf y}\le Ce^{-\gamma t}\abs{\mathbf y}
          \qquad\text{for all } t\ge0 \text{ and } \mathbf y\in\R^2 .
        \end{equation*}
        If $\mu_+\neq\mu_-$, one may also take $\gamma=\gamma_0$.
  \item If $\re\mu>0$ for an eigenvalue $\mu$ of $B$, then $x'=Bx$ has a real
        solution with $\abs{x(t)}\to\infty$ as $t\to\infty$.
\end{enumerate}
\end{proposition}

\begin{proof}
(a) \emph{Distinct eigenvalues.} Let $\bxi_\pm\in\C^2$ be eigenvectors for
$\mu_\pm$. They are linearly independent: $\bxi_-=c\,\bxi_+$ would give
$\mu_-\bxi_-=B\bxi_-=c\,\mu_+\bxi_+=\mu_+\bxi_-$. So the matrix $S$ with columns
$\bxi_+,\bxi_-$ is invertible, and $BS=S\Lambda$ with $\Lambda=\diag(\mu_+,\mu_-)$.
Then $B^k=S\Lambda^kS^{-1}$ for all $k$, and summing the exponential series gives
$e^{tB}=Se^{t\Lambda}S^{-1}$ with $e^{t\Lambda}=\diag(e^{\mu_+t},e^{\mu_-t})$. Since
$\abs{e^{\mu t}}=e^{t\re\mu}$, for $\mathbf y\in\R^2$
\begin{equation*}
  \abs{e^{tB}\mathbf y}\le\norm{S}\,\norm{e^{t\Lambda}}\,\norm{S^{-1}}\,\abs{\mathbf y}
  = Ce^{-\gamma_0t}\abs{\mathbf y}\le Ce^{-\gamma t}\abs{\mathbf y},
  \qquad C=\norm{S}\,\norm{S^{-1}},
\end{equation*}
and $C\ge1$ because $1=\abs{SS^{-1}\mathbf z}\le\norm{S}\norm{S^{-1}}$ for
$\abs{\mathbf z}=1$.

\emph{Repeated eigenvalue.} Now $\mu_\pm=\mu=\frac12\tr B$ is real and
$\mu^2=\det B$. Multiplying out shows $B^2-(\tr B)B+(\det B)I=0$ (the
Cayley--Hamilton theorem), so $N=B-\mu I$ satisfies
$N^2=B^2-2\mu B+\mu^2I=0$. Since $\mu I$ and $N$ commute,
$e^{tB}=e^{\mu t}e^{tN}=e^{\mu t}(I+tN)$; the series for $e^{tN}$ breaks off after
two terms. Given $\gamma<\gamma_0=-\mu$, let $\eps=\gamma_0-\gamma>0$. Then
\begin{equation*}
  \abs{e^{tB}\mathbf y}\le e^{\mu t}\big(1+t\norm{N}\big)\abs{\mathbf y}
  = e^{-\gamma t}\big(1+t\norm{N}\big)e^{-\eps t}\abs{\mathbf y}.
\end{equation*}
The function $t\mapsto(1+t\norm{N})e^{-\eps t}$ is continuous on $[0,\infty)$,
equals $1$ at $t=0$ and tends to $0$ as $t\to\infty$, so it is bounded by some
$C\ge1$.

(b) Let $\mu=\alpha+i\beta$ with $\alpha>0$, and let $\bxi$ be an eigenvector.
Then $e^{\mu t}\bxi$ solves $x'=Bx$, and since $B$ is real, so does its real part
$x(t)=e^{\alpha t}\big(\cos(\beta t)\re\bxi-\sin(\beta t)\im\bxi\big)$. If
$\beta=0$, we may take $\bxi$ real, and $\abs{x(t)}=e^{\alpha t}\abs{\bxi}\to\infty$.
If $\beta\neq0$, then $\re\bxi$ and $\im\bxi$ are linearly independent: otherwise
$\bxi$ would be a complex multiple of a real vector $\mathbf r\neq0$, and
$B\mathbf r=\mu\mathbf r$ with $B\mathbf r$ real would force $\mu\in\R$. Hence the
continuous $2\pi$-periodic function $s\mapsto\abs{\cos(s)\re\bxi-\sin(s)\im\bxi}$
has no zero, its minimum $c$ is positive, and $\abs{x(t)}\ge c\,e^{\alpha t}\to\infty$.
This argument works for real $n\times n$ matrices as well.
\end{proof}
